%HOMEWORK template for M 325K
\documentclass{article}
\headheight=7pt         \topmargin=14pt
\textheight=574pt       \textwidth=445pt
\oddsidemargin=18pt     \evensidemargin=18pt

%Begin package import

\usepackage{amsmath, amssymb, amsthm}
\usepackage{mathtools}
\usepackage{pdfpages}
\usepackage{tikz-cd}

%End package import
%Begin custom commands

\newcommand{\C}{\mathbb{C}}
\newcommand{\R}{\mathbb{R}}
\newcommand{\Q}{\mathbb{Q}}
\newcommand{\Z}{\mathbb{Z}}
\newcommand{\N}{\mathbb{N}}
\newcommand{\abs}[1]{\lvert #1\rvert}
\newcommand{\RM}[1]{\mathrm{#1}}
\newcommand{\BF}[1]{\textbf{#1}}
\newcommand{\e}{\epsilon}
\newcommand{\ceil}[1]{\lceil{#1}\rceil}
\newcommand{\floor}[1]{\lfloor{#1}\rfloor}
\newcommand{\rarr}[0]{\rightarrow}
\newcommand{\IM}[1]{\text{Im}(#1)}
\newcommand{\Hom}[0]{\text{Hom}}
\newcommand{\Pow}[1]{\text{Pow}(#1)}
\newcommand{\angles}[1]{\langle #1 \rangle}

%End custom commands
%Begin custom theorems

\newtheorem{innercustomgeneric}{\customgenericname}
\providecommand{\customgenericname}{}
\newcommand{\newcustomtheorem}[2]{%
\newenvironment{#1}[1]
{%
\renewcommand\customgenericname{#2}%
\renewcommand\theinnercustomgeneric{##1}%
\innercustomgeneric%
}
{\endinnercustomgeneric}
}

\newcustomtheorem{THRM}{Theorem}
\newcustomtheorem{LEM}{Lemma}
\newcustomtheorem{EX}{Exercise}
\newcustomtheorem{COR}{Corollary}
\newcustomtheorem{PROB}{Problem}
\newcustomtheorem{claim}{Claim}

\newenvironment{solution}
{\renewcommand\qedsymbol{$\blacksquare$}\begin{proof}[Solution]}
{\end{proof}}

\newenvironment{definition}
{\renewcommand\qedsymbol{$\blacksquare$}\begin{proof}[Definition]}
{\end{proof}}

%End custom theorems
\begin{document}

\title{Irreducible Reps of $D_n$}
\author{Arham Lodha}%put your name here
\date{January 31, 2024}%enter the due date here
\maketitle

\begin{PROB}{1}\label{Problem 1.}
    Let A and B be linear transformations of a finite dimensional vector space V. Show that if AB = BA then A and B preserve each others eigenspaces. If A is diagonalizable show the converse holds, ie a linear transformation B will commute with A iff B preserves the eigenspaces of A. Let G be a finite Abelian group. Show that every irreducible complex representation is one dimensional
\end{PROB}
\begin{claim}{1a}\label{Claim 1a.}
    Let A and B be linear transformations of a finite dimensional vector space V. Show that if AB = BA then A and B preserve each others eigenspaces.
\end{claim}
\begin{proof}
    For a linear transformation $T$, let $V_\lambda(T)$ be the eigenspace associated with $\lambda$.   

    Let $L_A$ be the set of $A$'s eigenvalues. Let $L_B$ be the set of $B$'s eigenvalues. $\forall \lambda \in L_A$, $\forall v \in V_\lambda(A)$, $A(Bv) = BAv = B \left(\lambda v\right) = \lambda Bv \implies Bv \in V_\lambda(A)$. Thus $B$ preserves $A$'s eigenspaces. $\forall \lambda \in L_B$, $\forall v \in V_\lambda(B)$, $B(Av) = ABv = A \left(\lambda v\right) = \lambda Av \implies Av \in V_\lambda(B)$. Thus $A$ preserves $B$'s eigenspaces.
\end{proof}
\begin{claim}{1b}\label{Claim 1b.}
    Suppose A is diagonalizable and B preserves the eigenspaces of A, then B will commute with A.
\end{claim}
\begin{proof}
    Suppose A is diagonalizable and B preserves the eigenspaces of A. From Linear Bootcamp, since A is diagonalizable there exists a basis eigenvalues of A of $v_1, \dots, v_k$. Let $\lambda_i$ be the associated eigenvalue to $v_i$. $\forall w \in V$, $w = \sum_{i = 1}^{k}a_iv_i$, thus $ABw = AB(\sum_{i = 1}^{k}a_iv_i) = A(\sum_{i = 1}^{k}a_iBv_i)$. Since $B$ preserves the eigenspaces of $A$, $Bv_i  \in V_{\lambda_i}(A)$. Thus $A(\sum_{i = 1}^{k}a_iBv_i) = \sum_{i = 1}^{k}a_i\lambda_iBv_i = \sum_{i = 1}^{k}a_iB\lambda_iv_i = \sum_{i = 1}^{k}a_iBAv_i = BA\sum_{i = 1}^{k}a_iv_i = BAw$. Thus $AB = BA$.      
\end{proof}
\begin{claim}{1c}\label{Claim 1c.}
    Let G be a finite Abelian group. Show that every irreducible complex representation is one dimensional.
\end{claim}
\begin{proof}
    Suppose $V$ is an irreducible complex representation. Let $f: G \to GL(V)$ be the representation. Assume the contrary, that $\dim V > 1$. Since $G$ is finite Abelian, all elements of g commute and thus by claim 1a, they preserve each others eigenspace. If $A \in GL(V)$, $\exists \lambda \in F$ $\ni V_\lambda(A)$. By the linear algebra bootcamp, we know that the only subspace of V with dimension equal to $\dim V$ is $V$. Thus $\forall v \in V$, $Av = \lambda v$. Thus, A is scalar. Suppose $\forall g \in G$, $f(g) = \lambda_g I$. Then $\forall v \in V \ni v \neq 0$, $\forall g \in G $, $f(g)v = \lambda_gv \in span(v)$. So $span(v)$ is G-invariant. By Mashke's Theorem and the previous homework, $\exists W$ a g invariant subspace complement to $span(v)$ thus V can be broken down into 2 complementary g-invariant subpaces thus V is reducible. Thus by contradiction, V must have dimension 1 if all $\forall g \in G$, $f(g)$ has a eigenspace of dimension equal to the dimension of the vectorspace. 
    
    Suppose $\exists g \in G \ni dim V_{\lambda}(f(g))< \dim V$ where $\lambda$ is a eigenvalue of $f(g)$. Since $G$ is finite Abelian, all elements of g commute and thus by claim 1a, they preserve each others eigenspace. So all elements of $G$ preserve $V_{\lambda}(f(g))$ thus $V_{\lambda}(f(g))$ is G-invariant. By Mashke's Theorem and the previous homework, $\exists W$ a g invariant subspace complement to $span(v)$ thus V can be broken down into 2 complementary g-invariant subpaces thus V is reducible. Thus by contradiction, V must have dimension 1. 
    
    Thus $V$ must have dimension 1. 
\end{proof}

\bigskip

\begin{claim}{2}\label{Claim 2.}
    $D_n$ "is" the group given by generators x,y with relations $x^n=y^2=1$, $yxy= x^{-1}$ What does this mean? State the universal property for a group given by some (specific) generators and relations.
\end{claim}
\begin{proof}
    Let $S = \left\{ x, y \right\}$. Let $G$ be a group given by the generators $x, y$ and relations $x^n = y^2 = 1$ and $yxy^{-1} = x^{-1}$. Let $\alpha: S \to G$ where $x \to x$ and $y \to y$. Let $\beta: S \to D_n$ where $x \to x$ and $y \to y$. By the universal property, $\exists$ unique homomorphisms $A: D_n \to G$ and $B: G \to D_n$ such that $\alpha = A \circ \beta$ and $\beta = B \circ \alpha$. Since $A$ and $B$ are unique homomorphisms, $A \circ B$ and $B \circ A$ must be unique homomorphisms. $A \circ B $ is the unique homomorphisms from $D_n \to D_n$, the identity is a homomorphism from $D_n \to D_n$. Since $A \circ B$ is the unique homomorphism from $D_n \to D_n$ $A \circ B = id_G$. Similarly since $B \circ A$ are the unique homomorphism from G to G, but the identity exists so $B \circ A = id_{D_n}$. Thus $A$ and $B$ are bijections and isomorphisms. Thus $G \cong D_n$. Thus up to isomorphisms, $D_n$ is the "is" the group given by generators x,y with relations $x^n=y^2=1$, $yxy^{-1} = x^{-1}$.                            
\end{proof}

\begin{claim}{3a}\label{Claim 3a.}
    Let V be an irreducible finite dimensional representation. Explain why x and y are diagonalizable (see the remarks below about a basic diagonalizability result).
\end{claim}
\begin{proof}
    Let $f: D_n \to GL(V)$ be a representation. 
    Let V be an irreducible finite dimensional.

    Since \[
        I = f(e) = f(x^n) = f(x)^n.
    \] and \[
        I = f(e) = f(y^2) = f(y)^2
    \].

    Thus $f(x)$ satsifies the polynomial $p_x(t) = t^n - 1$ and $f(y)$ satisfies the polynomial $p_y(t) = t^2 - 1$. The 2nd and nth roots of unitys are the roots of $p_y(t)$ and $p_x(t)$ and thus are distinct. By the linear algebra bootcamp and the remarks below, $f(x)$ and $f(y)$ are diagonalizable and the eigenvalues are the nth roots of unity and the 2nd roots of unity respectively     
\end{proof}
\begin{claim}{3b}\label{Claim 3b.}
    Suppose $f(x)v = \lambda v$ .
    Show that $f(x) (f(y)v) = (\bar \lambda) f(y) v$.  
\end{claim}
\begin{proof}
    $v = f(x)^{-1}f(x)v = f^{-1}(x) \lambda v = \lambda f(x)^{-1} v \implies \frac{1}{\lambda}v = f^{-1}(x)v$. $\frac{1}{\lambda} = \bar{\lambda}$, since $\lambda$ is a nth root of unity.   

    Since $yxy = x^{-1} \implies xy = y^{-1}x^{-1} = yx^{-1}$, 

    \[f(x)f(y)v = f(xy)v = f(yx^{-1})v = f(y)f(x)^{-1}v = f(y)\bar{\lambda}v\]
\end{proof}
\begin{claim}{3c}\label{Claim 3c.}
    Show that $span{v, f(y) v}$ is $D_n$ invariant.
\end{claim}
\begin{proof}
    $\forall w \in span(v, f(y)v)$, $w = av+ bf(y)v$ for $a, b \in F$. $f(x)w = f(x)(av + bf(y)v) = af(x)v + bf(x)f(y)v = a\lambda v + f(x)f(y)v = a \lambda v + \frac{1}{\lambda} f(y)v \in span(v, f(y)v)$. Thus $f(x)$ preserves $span(v, f(y)v)$. Suppose $f(x)^k w = w' \in span(v, f(y)v)$, then $f(x)^{k+1}w = f(x)f(x)^{k}w = f(x)w' \in span(v, f(y)v)$, since $f(x)$ preserves $span(v, f(y)v)$. Thus $\forall k \in \N$, $f(x^k) = f(x)^k$ preserves $span(v, f(y)v)$.
    
    $f(y) * v =f(y)v \in span(v, f(y)v)$ and $f(y) f(y)v = f(y^2)v = v \in span(v, f(y)v)$. Thus $f(y)$ swaps $v$ and $f(y)v$. Thus $f(y)$ preserves $span(v, f(y)v)$.        
    
    $\forall g \in D_n$, from last semester we know, $g = x^ay^b$ where $a = 1, \dots, n$ and $b = 0, 1$. Since $f(x^a)$ and $f(y^b)$ ($f(y^0) = I$) preserve $span(v, f(y)v)$. Since $f(x^a)f(y^b) = f(x^ay^b) = f(g)$ $f(g)$ preserves $span(v,f(y)v)$. 
    
    Thus $span(v, f(y)v)$ is $D_n$ -invariant. 
\end{proof}

\begin{claim}{3d}\label{Claim 3d.}
    Conclude that if V is irreducible, it has dimension either 1 or 2.
\end{claim}
\begin{proof}
    Suppose V is irreducible. By definition, there are no proper subspaces which are G-invariant. By the previous problem, $span(v, f(y)v)$ is G-invariant. Thus $span(v, f(y)v) = V$. $\dim span(v, f(y)v) = 1, 2$.      
\end{proof}


\begin{claim}{3e}\label{Claim 3e.}
    Suppose the dimension is 1. Show $x^2$ is in the kernel
\end{claim}
\begin{proof}
    Suppose $\dim V = 1$. Thus $\forall A \in GL(V)$, A is are scalar transformations. $\forall g \in G_n$ let $f(g) = \lambda_g I$. Since $f(y)f(y^{-1}) = f(y)f(y) = I \implies \lambda_y \lambda_{y^{-1}} I = \lambda_y^2 I = \lambda_y I \implies \lambda_y^2 = 1$. Since $\lambda_y  I= f(y) = f(y^{-1}) = \lambda_y^{-1} I$.  
    
    Since $\lambda_x I = \lambda_y \lambda_x \lambda_y I = f(y)f(x)f(y)^{-1} = f(x^{-1}) = \lambda_{x^{-1}} I \implies \lambda_x = \lambda_{x^{-1}}$. Since $f(x)f(x^{-1}) = \lambda_x \lambda_{x^{-1}} I = f(e) = I$. Thus $ I = \lambda_x \lambda_{x^{-1}} I = \lambda_{x} \lambda_{x} I = f(x)f(x) = f(x^2)$. Thus $x^2 \in \ker f$.   
\end{proof}

\begin{claim}{3f}\label{Claim 3f.}
    Classify all the irreducible reps and check that there are the right number (i.e. the number of conjugacy classes).
\end{claim}
\begin{proof}
    If $n$ is even, reflections through vertices are conjugate, reflections through edges are conjugate, and rotations are conjugate to their inverse. Since rotations are conjugate to their inverse $x^{\frac{n}{2}}$ is conjugate to itself. So the number of conjugacy classes is 2(size of the center) + $\frac{n - 2}{2}$(number of noncentral rotations) + 1 (vertex reflections) + 1(edge reflections) = $\frac{n - 2}{2} + 4$.
    
    If $n$ is odd, all reflections are conjugate, and rotations are conjugate to their inverse. So the number of conjugacy classes is 1(size of the center) + $\frac{n - 1}{2}$(number of noncentral rotations) + 1 (reflections) = $\frac{n - 1}{2} + 2$. 

    \textbf{Dimension 1}: From the last claim, since 
    
    \[\lambda_x^2 I = f(x)^2 = f(x^2) = I.\]

    This means $\lambda_x = \pm 1$ (can equal). Similarly since, 
    
    \[\lambda_y^2 I = f(y)^2 = f(y^2) = I.\]

    This means, $\lambda_y = \pm 1$. However if $n$ is odd, then $f(x)^n = \lambda_x^n I = I  \implies \lambda_x^n = 1$. Assume the contrary, $\lambda_x = -1$, however, $\lambda_x^n = (-1)^n = -1$, since $n$ is odd. Thus $\lambda_x = 1$ if $n$ is even. If $n$ is even, if $\lambda_x = -1$, then $\lambda_x^n = 1$, so it does satify the necessary conditions. 
    
    Thus if $n$ is even, the possible reps send $x \to I$ and $y \to I$ or $y \to -I$. So 2 possible reps. If $n$ is even, the possible reps send $x \to I$ or $x \to -I$ and $y \to I$ or $y \to -I$. So 4 possible reps.
    
    \textbf{Dimension 2}: From claim 3b, we know, for a eigenvalue $\lambda$ of $f(x)$ and a eigenvector $v$, $f(x)f(y)v = \bar{\lambda} f(y)v \implies f(y)v \in V_{\bar{\lambda}}(f(x))$. Thus $f(x)$ has eigenvalues $\left\{ \lambda, \bar{\lambda} \right\}$.
    
    (Note from now on, $X$ means representation of $f(x)$ and similarly for $Y := f(y)$)

    Suppose $\lambda \in \R$, then $\lambda = \bar{\lambda}$. Then $v, Yv \in V_\lambda(X)$. $X$ preserves $V_\lambda(X)$. From claim 3c, Y swaps $v$ and $Yv$, thus it preserves, $V_\lambda(X)$. Thus $V_\lambda(X)$ is G-invariant, thus a contradiction occurs and $\lambda \not \in \R$.
    
    Thus $v, Yv$ are linearly independent as they are in different eigenspaces. Take $v, Yv$ as the basis. 

    Thus $X = \begin{bmatrix}
        \lambda & 0 \\ 0 & \bar{\lambda}
    \end{bmatrix}$ since $v \in V_\lambda(X)$ and $Yv \in V_{\bar{\lambda}}(X)$. By claim 3c, $Y = \begin{bmatrix}
        0 & 1 \\ 1 & 0
    \end{bmatrix}$ since Y swaps $v$ and $Yv$.

    Suppose $\omega = \bar{\lambda}$, and say $X = \begin{bmatrix}
        x & 0 \\ 0 & \bar{x}
    \end{bmatrix}$ then we have simply swapped $v$ and $Yv$ because $v \in V_{\bar{x}}(X) = V_{\lambda}(X)$ and $Yv \in V_{x}(X) = V_{\bar{\lambda}}(X)$. If two representations are the same under change of basis, they are isomorphic. 

    Morever, representations with different characters are not isomorphic by the contrapositive of the converse of Frobenius Theorem (proved last homework). $tr(X) = \lambda + \bar{\lambda}$
    
    Let $\omega = e^{\frac{2\pi i}{n}}$.

    Thus if n is even, the possible pairs of eigenvalues of x are $\left\{ \omega, \omega^{-1} \right\}, \dots, \left\{ \omega^{\frac{n - 2}{2}}, \omega^{-\frac{n - 2}{2}} \right\}$. Thus if $n$ is odd, the possible pairs of eigenvalues are, $\left\{ \omega, \omega^{-1} \right\}, \dots, \left\{ \omega^{\frac{n - 1}{2}}, \omega^{-\frac{n - 1}{2}} \right\}$.

    



    % By 3a, we know the eigenvalues of $y$ are 1 and $-1$ and the eigenvalues of x part of the nth roots of unity and the its two eigenvalues are conjugates (by the above). 

    % Assume the contrary, x has real eigenvalues, then $\lambda = \bar{\lambda}$ and x is scalar, and thus preserves any subspace. However since $y$ is diagonalizable, the sum of the dimensions of the eigenspaces is $\dim V$. Moreover from 3a, the eigenvalues of $y$ can be 1 or -1. Suppose $y = \pm I$, then both $x$ and $y$ preserve any proper subspace of $V$ since they are both scalar and thus the subpace is G-invariant. By Mashke's theorem, $V$ is reducible. Thus by contradiction, $y \neq \pm I$. Suppose y's eigenvalues are $1$ and $-1$. Then $x$ preserves both the $V_1(y)$ subspace and $V_{-1}(y)$ subspace, thus both subspaces are G-invariant and complementary. Thus $V$ is reducible to $V_1(y) \times V_{-1}(y)$ which is a contradiction.        
        
    % Assume the contrary, $y = \pm I$. Then y preserves any subspace. Since x is diagonalizable, whose eigenvalues are part of the nth roots of unity. By the above $x$ has no real eigenvalues and thus isn't scalar. Let the eigenvalues of $x$ be $\omega$ and $\bar \omega$. Then $y$ preserves both the $V_\omega(x)$ and $V_{\bar{\omega}}(x)$ eigenspaces of X. Thus the two eigenspaces are G-invariant. By Mashke's Theorem, $V$ is reducible, which is a contradiction. Thus $y \neq \pm I$ and the eigenvalues of y are 1 and -1. 

    
    % Assume the contrary, the eigenvalue of x is 1, then its other eigenvalue must also be 1 because if $\lambda$ is a eigenvalue $\bar{\lambda}$ is also a eigenvalue, as stated above. Thus 1 is not a eigenvalue of x.
    
    % Assume the contrary, the eigenvalue of x is -1, then its other eigenvalue must also be 1 because if $\lambda$ is a eigenvalue $\bar{\lambda}$ is also a eigenvalue, as stated above. Thus -1 is not a eigenvalue of x.
    
    % Thus if n is even, the possible pairs of eigenvalues of x are $\left\{ \omega, \omega^{-1} \right\}, \dots, \left\{ \omega^{\frac{n - 2}{2}}, \omega^{-\frac{n - 2}{2}} \right\}$. Thus if $n$ is odd, the possible pairs of eigenvalues are, $\left\{ \omega, \omega^{-1} \right\}, \dots, \left\{ \omega^{\frac{n - 1}{2}}, \omega^{-\frac{n - 1}{2}} \right\}$.

    % Note in the following, I write $x \to A$ and $y \to B$ for some $A, B \in GL_2(\C)$, what I mean is that under the representation $x$ is sent to $PAP^{-1}$ and $y$ is sent to $PBP^{-1}$ for any $P \in GL_2(\C)$. Basically up to isomorphisms of $\C^2$ , $x \to A$ and $y \to B$ are the representations.
    
    $\forall k \in \left\{ 1, \floor{\frac{n - 1}{2}} \right\}$.

    $YXY = \begin{bmatrix}
        0 & 1 \\ 1 & 0
    \end{bmatrix}\begin{bmatrix}
        \omega^k & \\
        & \omega^{-k}
    \end{bmatrix}\begin{bmatrix}
        0 & 1 \\ 1 & 0
    \end{bmatrix} = \begin{bmatrix}
        0 & 1 \\ 1 & 0
    \end{bmatrix}\begin{bmatrix}
        & \omega^k \\
        \omega^{-k} & 
    \end{bmatrix} = \begin{bmatrix}
        \omega^{-k} & \\ & \omega^k
    \end{bmatrix} = X^{-1} = \begin{bmatrix}
        \omega^{-k} & \\ & \omega^k
    \end{bmatrix}$. Thus this satisfies the condition, $yxy = x^{-1}$, and thus for any k $X = \begin{bmatrix}
        \omega^k & \\
        & \omega^{-k}
    \end{bmatrix}$ and $Y = \begin{bmatrix}
        0 & 1 \\ 1 & 0
    \end{bmatrix}$ is a valid representation. 
    
    Thus the 2D irreducible representation of $D^n$ send $x \to \begin{bmatrix}
        \omega^k & \\
        & \omega^{-k}
    \end{bmatrix}$ for $k \in \left\{ 1, \floor{\frac{n - 1}{2}} \right\}$ and $y \to \begin{bmatrix}
        0 & 1 \\ 1 & 0
    \end{bmatrix}$. For $n$ odd, the number of 2D irreducible representations is $\frac{n - 1}{2}$. For $n$ even, the number of 2D irreducible representations is $\frac{n - 2}{2}$. 
    
    The total number of irreducible representations:

    For n even, there are 4 1D irreducible representations and $\frac{n - 2}{2}$ 2D irreducible representations. Thus $\frac{n - 2}{2} + 4$ irreducible representation. 
    
    For n odd, there are 2 1D irreducible representation and $\frac{n - 1}{2}$ 2D irreducible representation, thus $\frac{n - 1}{2} + 2$ irreducible representation.
    
    This matches the number of conjugacy classes thus we have found all of them. 

\end{proof}

\bigskip

\begin{claim}{4}\label{Claim 4.}
    Now let G be a non-abelian group of order 21. Show $G = \angles{x,y| y^3=x^7=1, y^{-1}xy = x^2}$ .
\end{claim}
\begin{proof}
    Let $G$ be a non-abelian group where $|G| = 21 = 3 * 7$. By the Sylow theorem I, $\exists H \leq G \ni |H| = 7$. Since $7$ is prime, $H$ must be cyclic. By Sylow Theorem III, the number of conjugate subgroups of $H$ is $[N(H) : H] = |N(H) / H|$ and $[N(H) : H]$ divides 3 and $[N(H) : H] \mod 7 = 1$. Since 3 is prime, the only numbers which divide it are 1 and 3. Assume the contrary, $[N(H) : H] = 3$ but $[N(H) : H] \mod 7 = 3$ and this leads to a contradiction. Thus $[N(H) : H]=1$. Thus there is only 1 conjugate subgroup of order 7. Thus $H$ is normal. Since $H$ is cyclic order 7, $H = C_7$ (Going forward $C_7$ is used). By the Sylow Theorem 1, $\exists A \leq G \ni |A| =3$. Since $3$ is prime, $A = C_3$. 

    $C_3 \cap C_7 \leq C_3$ and $C_3 \cap C_7 \leq C_7$ so $|C_3 \cap C_7| | 3$ and $|C_3 \cap C_7| | 7$ thus $|C_3 \cap C_7 |= 1$, thus $C_3 \cap C_7 = 1$.     
    
    By direct product (last semester), $|C_3 \times C_7| = |C_3||C_7|/\gcd(7, 3) = |C_3||C_7|=21$. $\forall (a, b), (c,d) \in C_3 \times C_7, (a, b) * (c, d) = (ac, bd)$. Since $C_3$ and $C_7$ are abelian, $ac = ca$ and $bd= db$. Thus $(a, b) * (c, d) = (ac, bd) = (ca, db) = (c,d)*(a, b)$. Thus $C_3 \times C_7$. Thus $G \not \cong C_3 \times C_7$. 

    Note $y$ is used as generator of $C_3$ and $x$ used as generator of $C_7$, and $C_7$ .    
    
    $\forall \phi \in Aut(C_7)$, takes $g \to g^k$ where $k$ is relatively prime to 7.  Thus $Aut(C_6) \to \left\{ 1, 2, 3, 4,5,6 \right\}$. The homomorphism $\phi: C_3 \to Aut(C_7)$ is nontrivial. Assume the contrary $\phi$ is trivial map, $C_7 \rtimes_\phi C_3 \cong \angles{x, y | x^7= y^3 = e, yxy^{-1} = x} = \angles{x, y | x^7= y^3 = e, yx = xy}$, which is abelian. 
    
    Thus $\phi$ is nontrivial, Thus $\phi: C_3 \to C_6$ takes $x \to y^k$ for some $k \in {0, \dots, 5}$ . $e = \phi(y^3) = \phi(y)^3 = (y^k)^3 = x^{3k} \implies 3k \equiv 0 \mod 6 \implies k = 2, 4$. 
    
    Suppose $\phi_1(y) = x^2$. Thus $C_7 \rtimes_{\phi_1} C_3 = \angles{x, y | x^7 = y^3 = e, y^{-1}xy = x^2}$. 
    
    Suppose $\phi_2(x) = y^4$. Thus $C_7 \rtimes_{\phi_2} C_3 = \angles{x, y | x^7 = y^3 = e, y^{-1}xy = x^4}$. Let $y^{-1}$ also generates $C_3$, so let $b = y^{-1}$. Let $x^4$ also generates $C_7$, so let $a = x^4$. thus $yxy^{-1} = x^4 \implies x = yx^{-4}y^{-1} \implies a^2 = b^{-1}ab$. Thus $C_7 \rtimes_\phi C_3 \cong \angles{a, b | a^7 = b^3 = 1, b^{-1}ab = a^2}$. Thus $C_7 \rtimes_{\phi_2} C_3 \cong C_7 \rtimes_{\phi_1} C_3$          
    
    Thus the only group of order 21 is $\angles{x, y | x^7 = y^3 = e, y^{-1}xy = x^2} = G$.      
\end{proof}

\bigskip

\begin{PROB}{5}\label{Problem 5.}
    $G = \angles{x,y| y^3=x^7=1, y^{-1}xy = x^2}$
    Show every irreducible representation of G has dim 1 or 2, and classify the irreducible reps.
\end{PROB}
\begin{PROB}{5a}\label{Problem 5a.}
    $x$ and $y$ are diagonalizable  
\end{PROB}
\begin{proof}
    Let $f: G \to GL(V)$ be a representation. 
    Let V be an irreducible finite dimensional.

    Since \[
        I = f(e) = f(x^7) = f(x)^7.
    \] and \[
        I = f(e) = f(y^3) = f(y)^3
    \].

    Thus $f(x)$ satsifies the polynomial $p_x(t) = t^7 - 1$ and $f(y)$ satisfies the polynomial $p_y(t) = t^3 - 1$.
    
    The 3nd and 7th roots of unitys are the roots of $p_y(t)$ and $p_x(t)$ and thus are distinct.
    
    Thus by the linear algebra bootcamp and the remarks below, $f(x)$ and $f(y)$ are diagonalizable. The eigenvalues of $f(x)$ and $f(y)$ are the 7th and 3rd roots of unity respectively.   
\end{proof}
\begin{claim}{5b}\label{Claim 5b.}
    Suppose $f(x)v = \lambda v$ .
    Show that $f(x)f(y)v = \lambda^2 f(y)v$ and that $f(x)f(y^2) v = \lambda^4 f(y^2)v$ . 
\end{claim}
\begin{proof}  
    Since $y^{-1}xy = x^{2} \implies x^2y = yx$, 

    \[f(x)f(y)v = f(xy)v = f(y y^{-1} x y)v = f(yx^2)v = f(y)\lambda^2v = \lambda^2 f(y) v\]

    Thus $f(y)v \in V_{\lambda^2}(f(x))$. 
    
    
    \[
        f(x)f(y^2)v = f(y y^{-1} x y y)v = f(y x^2 y)v = f(y x x y)v = f(y x y x^2) v =f(y^2 x^4) v = \lambda^4 f(y^2) v
    \].

    Thus $f(y^2)v \in V_{\lambda^4}(f(x))$. 
\end{proof}

% \begin{claim}{5c}\label{Claim 5c.}
%     $\forall g \in G$, $g = x^ay^b$ where $a \in \left\{ 0, 1, \dots, 7 \right\}$ and $b = \left\{0, 1, 2 \right\}$
% \end{claim}
% \begin{proof}
         
% \end{proof}

\begin{claim}{5c}\label{Claim 5c.}
    Show that $span(v, f(y)v, f(y^2)v)$ is $G$-invariant.
\end{claim}
\begin{proof}
    $f(x)v = \lambda v \in span(v, f(y)v, f(y^2)v)$. By claim 5b, $f(x)f(y)v = \lambda^2 f(y) v \in span(v, f(y)v, f(y^2)v)$. By claim 5b, $f(x)f(y^2)v = \lambda^4 f(y^2) v \in span(v, f(y)v, f(y^2)v)$. Thus $f(x)$ preserves $span(v, f(y)v, f(y^2)v)$.
    
    $f(y)(v) = f(y)v \in span(v, f(y)v, f(y^2)v)$. $f(y)f(y)v = f(y^2)v \in span(v, f(y)v, f(y^2)v)$. $f(y)f(y^2)v = f(y^3)v = v \in span(v, f(y)v, f(y^2)v)$. Thus $f(y)$ preserves $span(v, f(y)v, f(y^2)v)$. Thus $f(y)$ sends $v \to f(y)v$, $f(y)v \to f(y^2)v$, and $f(y^2)v \to v$. Thus $f(y)$ preserves $span(v, f(y)v, f(y^2)v). $
    
    
    $f(y^2)(v) = f(y^2)v \in span(v, f(y)v, f(y^2)v)$. $f(y^2)f(y)v = v \in span(v, f(y)v, f(y^2)v)$. $f(y^2)f(y^2)v = f(y)v \in span(v, f(y)v, f(y^2)v)$. Thus $f(y^2)$ sends $v \to f(y^2)v$, $f(y)v \to v$, $f(y^2)v \to f(y)v$. Thus $f(y^2)$ preserves $span(v, f(y)v, f(y^2)v)$.

    Since $x, y$ generate G, if they preserve $span(v, f(y)v, f(y^2)v)$, $\forall g \in G$, $f(g)$ preserves $span(v, f(y)v, f(y^2)v)$ because $g$ can be written as a finite string of $x$ and $y$. 
    Thus $span(v, f(y)v, f(y^2)v)$ is G-invariant.      
\end{proof}

\begin{claim}{5d}\label{Claim 5d.}
    Vectorspace $V$ irreducible has dimension 1, 2, or 3 
\end{claim}
\begin{proof}
    Suppose $V$ is irreducible. By definition, there are no proper subspaces which are G-invariant. By claim 5c, $span(v, f(y)v, f(y^2)v)$ is G-invariant. Thus $V = span(v, f(y)v, f(y^2)v)$. Thus $\dim V = 1, 2, 3$.      
\end{proof}

\begin{claim}{5e}\label{Claim 5e.}
    If $V$ is a dimension 1 irreducible representation, then $x \in \ker f$.  
\end{claim}
\begin{proof}
    If $\dim V = 1$, then $f: G \to GL_1(\C) = C^*$ and thus $\forall g \in G$, $f(g)$ commutes. 
    
    Thus, $f(x) = f(y)^{-1}f(x)f(y) = f(x^2) \implies f(x) = f(x)^{2} \implies f(x) = I$. Thus $x \in \ker f$  
\end{proof}

\begin{claim}{5e}\label{Claim 5e.}
    Classify all the irreducible reps and check that there are the right number (i.e. the number of conjugacy classes).
\end{claim}
\begin{proof}
    \textbf{Dimension 1}: From the last claim, $f(x) = I$ and $f(y) = \lambda_y I$. By the relations, $I = f(y^3) = \lambda_y^3 I \implies \lambda_y^3 = 1$. Thus $\lambda_y = e^{\frac{2 \pi i k}{3}}$ for $k \in \left\{ 0, 1, 2 \right\}$. Thus $f(x) = I$ and $f(y) = e^{\frac{2 \pi i k}{3}} I $ for $k \in \left\{ 0, 1, 2 \right\}$    
    
    (\textbf{Going forward, for the sake of notation, $X = f(x)$ and $Y = f(y)$})

    \textbf{Dimension 2}: 
    
    By the linear algebra bootcamp, since $x$ is diagonalizable, sum of the dimensions of the eigenspaces is 2. Suppose $\lambda$ is a eigenvalue of $X$ then $\lambda^2$ and $\lambda^4$ are also eigenvalues of x by 5b. 
    
    Asssume the contrary, $\lambda = \lambda^2$. Since $\lambda$ is a 7th root of unity. $\lambda = 1$. Then, $X$ is scalar and $X = I$. Thus X preserves any subspace of V. $X$ commutes with Y since $X = I$. Then the representation is abelian and V cannot be a dimension 2 irreducible vectorspace by problem 1. Thus $\lambda \neq \lambda^2$. 
    
    
    By linear algebra bootcamp, the number of distinct eigenvalues cannot exceed the dimension of the vectorspace. Thus by contradiction, there is no irreducible representation of G with dimension 2.   
    
    
    
             
    \textbf{Dimension 3}: $\dim span(v, Yv, Y^2v)$ must have dimension 3, otherwise it is a proper G-invariant subspace of V, which would make V reducible, thus causing a contradiction. Thus $\dim span(v, Yv, Y^2v)$ and $v, Yv, Y^2v$ is a basis of V. 

    Take $v, Yv, Y^2v$ as the basis.
    
    Let $\omega = e^{\frac{2 \pi i}{7}}$. By claim 5a, the possible eigenvalues are $\lambda \in \left\{ \omega^k | k \in \left\{ 0, 1, 2, 3, 4, 5, 6 \right\} \right\}$. 
    
    Thus $\forall k \in \left\{ 0, \dots, 6 \right\}$, $X = \begin{bmatrix}
        \lambda \\  & \lambda^2 \\  & & \lambda^4
    \end{bmatrix}$. Y takes $v \to Yv$ and $Yv \to Y^2v$ and $Y^2v \to v$. Thus $Y = \begin{bmatrix}
        0 & 0 & 1 \\ 1 & 0 & 0 \\  0 & 1 & 0
    \end{bmatrix}$.

    Assume the contrary, $k = 0$, $X = I$, then $X$ preserves any subspace of V. Thus X commutes with everything and the representation is abelian, thus by problem 1 the only irreducible representations are dimension 1. Thus by contradiction, $X \neq I$.
    
    Thus the possible values of $\lambda$ are $\left\{ \omega^k : k \in \left\{ 1, 2, 3, 4, 5, 6 \right\} \right\}$. 

    $\forall k \in \left\{1, 2, 3, 4, 5, 6 \right\}$
    
    $X^7 = \begin{bmatrix}
        \omega^k \\  & \omega^{2k} \\  & & \omega^{4k}
    \end{bmatrix}^7 = \begin{bmatrix}
        \omega^{7k} \\  & \omega^{7k} \\  & & \omega^{14k}
    \end{bmatrix} = I$.
    
    $Y^3 = \begin{bmatrix}
        0 & 0 & 1 \\ 1 & 0 & 0 \\  0 & 1 & 0
    \end{bmatrix}^3 = I
    $

    $Y^{-1}XY = \begin{bmatrix}
        0 & 0 & 1 \\ 1 & 0 & 0 \\  0 & 1 & 0
    \end{bmatrix}\begin{bmatrix}
        \omega^k \\  & \omega^{2k} \\  & & \omega^{4k}
    \end{bmatrix}\begin{bmatrix}
        0 & 0 & 1 \\ 1 & 0 & 0 \\  0 & 1 & 0
    \end{bmatrix}^{-1} = \begin{bmatrix}
        \omega^{2k} & 0 & 0 \\
        0 & \omega^{4k} & 0 \\
        0 & 0 & \omega^{k}
        \end{bmatrix} = X^2$
        
    Thus, all $k$ are valid.  

    By the contrapositive of the converse of the Frobenius theorem, which we proved last homework, if characters are nonequal then representations are not isomorphic. 
    
    Even if you reorder, $\omega^{2k}, \omega^{4k}, \omega^{8k}$ in X, you simply change the ordering of the basis and thus the representations are the same up to change of basis. 
    
    Thus representations where X has eigenvalues, $\omega^{2}, \omega^{2}, \omega^{4}$ are isomorphic and any representations where X has eigenvalues, $\omega^{3}, \omega^{6}, \omega^{5}$ are isomorphic. The representations where X has eigenvalues $\omega^{2}, \omega^{2}, \omega^{4}$ and the representations where $\omega^{3}, \omega^{6}, \omega^{5}$ are not isomorphic because their traces are not equal because $\omega^1 + \omega^2 + \omega^4 \neq \omega^3 + \omega^3 + \omega^5$.


    Thus the two representations are the following:

    \[
        x \to \begin{bmatrix}
            \omega \\  & \omega^{2} \\  & & \omega^{4}
        \end{bmatrix}, y \to \begin{bmatrix}
            0 & 0 & 1 \\ 1 & 0 & 0 \\  0 & 1 & 0
        \end{bmatrix},
    \]

    and

    \[
        x \to \begin{bmatrix}
            \omega^3 \\  & \omega^{5} \\  & & \omega^{6}
        \end{bmatrix}, y \to \begin{bmatrix}
            0 & 0 & 1 \\ 1 & 0 & 0 \\  0 & 1 & 0
        \end{bmatrix},
    \]

    Thus there are 2 irreducible representations of dimension 3, up to isomorphism. 

    Thus there are 5 representations of G in total. 

    
\end{proof}



\end{document} 